\documentclass[a4paper, 11pt]{article}

\usepackage[utf8]{inputenc}
\RequirePackage[left=2cm,text={17cm, 24cm},top=3cm]{geometry}
\usepackage{times}
\usepackage{graphicx}
\usepackage{setspace}
\usepackage{hyperref}
\usepackage{fancyhdr}
\usepackage[outputdir=build]{minted}
\usepackage{amsmath}
\usepackage{textcomp}

% for indented itemize list
\usepackage{enumitem}

\graphicspath{{../src/outputs/}}

\begin{document}

    % header
    \begin{center}
    {\includegraphics[width=0.4\linewidth]{fig/logo_en.png}}
        \\
        [0.4cm]

        {\LARGE
        ISS - Project \\
        [0.4cm]
        }

        {\large
        Adrián Kálazi (xkalaz00) \\
        \today
        }
    \end{center}

    % actual content
    \section*{Abstract}

    The calculations were implemented using \textbf{Python (3.5+)} with the following libraries:

    \begin{itemize}
        \item numpy - arrays, \texttt{fft}, \texttt{ifft}
        \item scipy - \texttt{lfilter}
        \item soundfile - wav reading/writing
        \item matplotlib - plotting
    \end{itemize}

    The implementation consists of a Python module and a launcher script (\texttt{main.py}).

    Running the script:
    \begin{itemize}
        \item
        On a system without the required libraries - execute \texttt{./run-venv.sh} in the \texttt{src/} folder.

        This will setup a virtual environment, install the requirements and run the script.
        \item
        On a system with the required libraries - execute \texttt{python3 ./main.py} in the \texttt{src/} folder.
    \end{itemize}

    \section*{Solution}

    \begin{enumerate}
        \item
        The audio was recorded with \texttt{\href{https://apps.kde.org/en/kwave}{KWave}} in 16-bit resolution at 48~kHz.

        The sampling rate was later scaled down to 16~kHz using \texttt{ffmpeg}

        Durations of recorded signals: \\
        \begin{tabular}{ | l | l | l | }
            \hline
            \textbf{File}     & \textbf{Samples} & \textbf{Seconds} \\ \hline
            maskoff\_tone.wav & 27436            & 1.71             \\ \hline
            maskon\_tone.wav  & 36696            & 2.29             \\ \hline
        \end{tabular}

        \item
        The sentence recording process was the same as above

        Durations of recorded signals: \\
        \begin{tabular}{ | l | l | l | }
            \hline
            \textbf{File}         & \textbf{Samples} & \textbf{Seconds} \\ \hline
            maskoff\_sentence.wav & 41162            & 2.57             \\ \hline
            maskon\_sentence.wav  & 43312            & 2.71             \\ \hline
        \end{tabular}

        \newpage
        \item
        Firstly, one second long signals were extracted from the wav files.

        Then the DC bias has been removed and the signals were normalized.

        The base tasks are calculated with 20ms frames extracted from the signal, resulting in a total of 99 frames
        (the last half-frame is ignored).

        We later recreate the frames with a 25ms duration for task 15.

        \vspace{2mm}
        \textbf{Calculation of the frame size:}

        Frame duration $t_F = 20\,ms$ \\
        Sampling frequency $F_s = 16\,kHz$

        Frame size (in samples): $n_F = t_F \times F_s = 20~ms \times 16~kHz = 320~samples$

        \includegraphics[width=\linewidth]{3_frames.pdf}

        \vspace{2mm}
        \item
        The correlation function implementation is located in \texttt{iss.operations.correlate\_frame}
        and the function is also used in task 15 for correlation.

        Mean and variance of base frequencies: \\
        \begin{tabular}{ | l | l | l | }
            \hline
            \textbf{Signal}       & \textbf{Mean} $[Hz]$ & \textbf{Variance} $[Hz]$ \\ \hline
            maskoff\_sentence.wav & 139.466              & 0.803                    \\ \hline
            maskon\_sentence.wav  & 143.478              & 1.883                    \\ \hline
        \end{tabular}

        \vspace{2mm}
        \textbf{Q:} How can the $f_0$ change be reduced for $\pm1$ errors of the lag?

        \textbf{A:} There are two possible solutions that come to my mind to the mentioned problem:

        \begin{itemize}
            \item
            In our implementation, lag was defined as \texttt{np.argmax(autocorrelation\_array)} and therefore it was an integer.

            An alternate implementation would be to approximate lag as a weighted arithmetic mean from the surrounding indices and their values
            which would make it a floating point number resulting in better precision.
            This method would eliminate rapid changes of $f_0$ in most cases.

            \item
            Autocorrelation with higher resolution - this would eliminate the rapid changes in $f_0$
            but would require a longer computation time.
        \end{itemize}

        \newpage
        \includegraphics[width=\linewidth]{4_frame.pdf} \\
        \includegraphics[width=\linewidth]{4_frame_clipped.pdf} \\
        \includegraphics[width=\linewidth]{4_frame_autocorrelated.pdf} \\
        \includegraphics[width=\linewidth]{4_base_frequencies.pdf}

        After center clipping and autocorrelation, the resulting base frequencies were examined
        and we came to the conclusion that the frequencies are "close enough"
        for them to be used safely for designing our filter.

        Task 11 also implements an alternative where only frames with matching base frequencies are used for
        determining the frequency response.

        \newpage
        \item
        For calculations, the fast fourier transform function (\texttt{fft} from \texttt{np.fft}) was used

        Example DFT function implementation

        % @formatter:off
        % implementation from https://www.fit.vutbr.cz/study/courses/ISS/public/NEW_PRED/02_spectrum/spectrum.pdf slide 46
        \begin{minted}[resetmargins=true,tabsize=4, python3=true]{python}
        def dft(x):
            X = []
            N = len(x)
            for k in range(N):
                a = 0
                for n in range(N):
                    a += x[n] * cmath.exp(-2j * cmath.pi * (k/N) * n)
                X.append(a)
            return X
        \end{minted}
        % @formatter:on

        \includegraphics[width=\linewidth]{5_spectrum_maskoff.pdf} \\
        \includegraphics[width=\linewidth]{5_spectrum_maskon.pdf}

        The spectrograms show a bigger presence of high-frequency noise in the "mask on" than the "mask off" signal

        \newpage
        \item
        Frequency response

        For the frequency response we need to calculate the the mean from all frames.

        \begin{gather*}
            H(e^{j\omega}) = \frac{1}{N} \sum_{n=0}^{N-1} \frac{ B_n(e^{j\omega}) }{ A_n(e^{j\omega}) } \\
            N \text{ - frame count, } A_n, B_n \text{ - spectrums for frame n}
        \end{gather*}

        \includegraphics[width=\linewidth]{6_frequency_response.pdf}

        From the filters frequency response we can see that the gain oscillates frequently
        over the lower half of the frequency range,
        which may be caused by the varying base frequency of the input signals.

        \item
        Impulse response

        The inverse fourier transform (\texttt{ifft} from \texttt{np.fft})
        was applied to the frequency response calculated in the previous step.

        \includegraphics[width=\linewidth]{7_impulse_response.pdf}

        IDFT function implementation

        % @formatter:off
        % implementation from https://www.fit.vutbr.cz/study/courses/ISS/public/NEW_PRED/02_spectrum/spectrum.pdf slide 51
        \begin{minted}[resetmargins=true,tabsize=4, python3=true]{python}
        def idft(X):
            x = []
            N = len(X)
            for n in range(N):
                a = 0
                for k in range(N):
                    a += X[k] * cmath.exp(2j * cmath.pi * (k/N) * n)
                x.append(a / N)
            return x
        \end{minted}
        % @formatter:on

        \newpage

        \item
        Filtered signal

        We applied the signal using \texttt{lfilter} from \texttt{scipy.signal}

        \includegraphics[width=\linewidth]{8_signal_maskoff.pdf} \\
        \includegraphics[width=\linewidth]{8_signal_filtered.pdf} \\
        \includegraphics[width=\linewidth]{8_signal_maskon.pdf}

        From the above figures we can see that the target signal and the filtered signal are different in most aspects.

        This is probably because of the differences in recording
        (we can see that the tones between the "mask off" and "mask on") recordings are not very similar.

        The filter seems to reduce a little bit of low frequency noise, but it's hardly noticeable.

        Overall the filter seems to "not do much",
        although in the recordings we can hear a little bit of a "dampening" effect.

        \item
        Conclusion

        The biggest problem we encountered was that we didn't have access to an example signal
        on which our calculations could be tested and this caused our results to be very dependent
        on our own recordings.

        This introduced quite a few problems
        as we didn't know if our calculations are right at any point.

        The filters weird behaviour may also be caused by a smaller difference between the tone recordings
        that the filter has been constructed from.

        Even though our calculations should be right, the results aren't good enough to be considered "acceptable".

    \end{enumerate}

    \newpage
    \section*{Solution - additional tasks}

    All of the next tasks recreate the signal frames,
    so that each task can be compared to the results from the base solution

    \begin{enumerate}
        \item[11.]
        Window function

        We created the window using \texttt{get\_window} from \texttt{scipy.signal}

        The slected window function that has been used was the "Hann window" with the following properties:

        \includegraphics[width=\linewidth]{11_window.pdf} \\
        \includegraphics[width=\linewidth]{11_response.pdf}

        Phase shift (from task 15) has been applied before the window function,
        this choice was made because it shows better how the function behaves in the next figures.

        \includegraphics[width=\linewidth]{11_frames_before.pdf} \\
        \includegraphics[width=\linewidth]{11_frames_after.pdf}

        \newpage
        Comparison of specral densities before and after applying the window function:

        \includegraphics[width=\linewidth]{11_frame_spectrum_before.pdf} \\
        \includegraphics[width=\linewidth]{11_frame_spectrum_after.pdf}

        From the above figures we can see that the window function
        supresses the frequencies on the edge of the spectrum and the resulting frequency spectrum is overall smoother.

        Window functions are useful for distributing the spectral leakage introduced by examining longer signals.

        \vspace{6mm}
        \item[13.]
        Filter constructed from only frames with matching base frequencies

        The needed frames were collected during task 4 and the frequency response
        (which is normally constructed in task 5) was instead contructed from these selected frames only.

        \includegraphics[width=\linewidth]{6_frequency_response.pdf} \\
        \includegraphics[width=\linewidth]{13_frequency_response.pdf}

        From the above figures, we can see that using only frames with matching base frequencies supresses
        sudden spikes in the frequency response

        \newpage
        \item[15.]
        Phase shift

        \includegraphics[width=\linewidth]{15_frames_before.pdf} \\
        \includegraphics[width=\linewidth]{15_frames_after.pdf} \\
        \includegraphics[width=\linewidth]{15_shifts.pdf}

        From the above graph, we can see that the required phase shift decreases almost linearly -
        this is caused by the frequency of the "mask on" signal being a little bit lower than the "mask off" signal

        In the best case scenario (when both signal frequencies match) the phase shift would be constant.

        \vspace{6mm}

        \textbf{Q:} When we take the sum of phase shifts from both sides,
        we can see that it resembles the lag in autocorrelation,
        what is the reason for this behaviour?

        \textbf{A:} This is due to the coefficients (and lag) representing a shift at which both signals are the "most similar".

        In autocorrelation this results in the lag representing the number of samples until the next period occurs.

        Because autocorrelation is done on only one signal,
        the shortest shift at which the signal is similar to itself is its own period.

        \textbf{The autocorrelation coefficient (and therefore lag) represents one period.}

        In cross-correlation the coefficient represents the shortest shift in a certain direction,
        but because the signals have similar frequencies and the cross-correlation is done from both sides,
        summing up these two values results in them
        being equal to the signals period $T$ in samples.

        \textbf{Both cross-correlation coefficients represent a part of the period (e.g.\ 1/3 and 2/3)
        which summed up results in one period.}
    \end{enumerate}

\end{document}
